% from ega0-23-fr.tex
\oldpage[0\textsubscript{IV-1}]{215}

raisonnement que ci-dessus montre que la fermeture intégrale $V'$ de $V$ dans
$K'$ est l'ensemble des $x'\in K'$ tels que $x'^q\in V$~; on sait
(Bourbaki, \emph{Alg. comm.}, chap.~VI, §~8, n\textsuperscript{o}~6,
prop.~6, et chap.~V, §~2, n\textsuperscript{o}~3, lemme~4) que $V'$ est un
\emph{anneau de valuation discrète}, dont l'idéal maximal $\mathfrak m'$ est
l'ensemble des $x'\in K'$ tels que $x'^q\in\mathfrak m$ et en outre le corps
résiduel $V'/\mathfrak m'$ est une extension \emph{finie} de
$V/\mathfrak m$ (\emph{loc. cit.}, chap.~VI, §~8,
n\textsuperscript{o}~1, lemme~2). Montrons que, pour tout entier $n>0$, on a
\begin{equation}
\label{0.23.1.3.2-fr}
  \mathfrak m'^n\cap A'=y^nA'.
  \tag{23.1.3.2}
\end{equation}

En effet, comme $y^q=x$, on a $y\in\mathfrak m'$, donc
$y^nA'\subset\mathfrak m'^n\cap A'$. Inversement, soit
$x'\in\mathfrak m'^n\cap A'$, et posons $x'=z'y^n$ avec $z'\in K'$. On a
$x'^q\in A'$~; d'autre part, $x'$ est somme de produits
$t'_1\cdots t'_n$ avec $t'_i\in\mathfrak m'$, donc
$t_i'^q\in\mathfrak m$, et l'on en conclut que
$x'^q\in\mathfrak m^n\cap A$. Or le lemme (23.1.3.1) prouve que
$\mathfrak m^n\cap A=x^nA$, donc on a $z'^qx^n\in x^nA$, ou encore
$z'^q\in A$ puisque $x\neq0$, ce qui établit (23.1.3.2).

Montrons en second lieu que sur $A'$ la topologie $xA'$-préadique est
\emph{séparée}~; cette topologie est aussi la topologie $yA'$-préadique de
$A'$, et il résulte de (23.1.3.2) que cette topologie est induite par la
topologie $\mathfrak m'$-préadique de $V'$ qui est séparée puisque $V'$ est
un anneau de valuation discrète.

Prouvons ensuite que $A'/xA'$ est un $A$-module de type fini. Comme $x=y^q$
et que $y^kA'/y^{k+1}A'$ est isomorphe à $A'/yA'$, on peut se borner à
montrer que $A'/yA'$ est un $A$-module de type fini~; mais la formule
(23.1.3.2) appliquée pour $n=1$ montre que $A'/yA'$ est un sous-anneau de
$V'/\mathfrak m'$, c'est-à-dire du corps résiduel de $V'$, qui est une
extension finie du corps résiduel $V/\mathfrak m$ de $V$. Or
$V/\mathfrak m$ est le corps des fractions de $A/\mathfrak p$, et comme
$A'$ est entier sur $A$, $A'/yA'$ est entier sur $A/\mathfrak p$, donc
contenu dans la fermeture intégrale de $A/\mathfrak p$ dans
$V'/\mathfrak m'$~; comme par hypothèse $A/\mathfrak p$ est un anneau
japonais noethérien, $A'/yA'$ est un $(A/\mathfrak p)$-module de type fini,
et \emph{a fortiori} un $A$-module de type fini.

Cela étant, la topologie $xA'$-préadique de $A'$ étant séparée, $A'$ est un
sous-anneau de son complété $\widehat{A'}$ pour cette topologie~; mais puisque
$A'/xA'$ est un $A$-module de type fini et que $A$ est séparé et complet pour
la topologie $xA$-préadique, $\widehat{A'}$ est un $A$-module de type fini en
vertu de (0\textsubscript{I}, 7.2.9), donc il en est de même de $A'$, C.Q.F.D.
\end{proof}

\begin{corollary}[23.1.4]
\label{0.23.1.4-fr}
Soit $A$ un anneau japonais noethérien et intégralement clos. Alors tout
anneau de séries formelles $A[[T_1,\ldots,T_r]]$ est un anneau japonais.
\end{corollary}

\begin{proof}
On sait que pour tout anneau noethérien intégralement clos $B$, l'anneau de
séries formelles $B[[T]]$ est noethérien et intégralement clos (Bourbaki,
\emph{Alg. comm.}, chap.~V, §~1, n\textsuperscript{o}~4, prop.~14)~; on peut
donc, par récurrence sur $n$, se borner à prouver que $A[[T]]$ est un anneau
japonais. Or l'élément $x=T$ vérifie toutes les conditions de (23.1.3), d'où
la conclusion.
\end{proof}

\begin{theorem}[23.1.5]
\label{0.23.1.5-fr}
\textnormal{(Nagata).} Tout anneau local noethérien intègre complet $A$ est
un anneau japonais.
\end{theorem}

\begin{proof}
On sait (19.8.8, (ii)) qu'il existe un sous-anneau $B$ de $A$ qui est un
anneau local complet et régulier, tel que $A$ soit une $B$-algèbre finie~; il
suffit évidemment de prouver

